\section{Content-verified retrieval}
\label{sec:retrieval}
\label{sec:lanes}

Every citation a lane proposes must pass a content-verification gate before it
may appear in an answer. The gate scores each retrieved segment on a 0--8 rubric
(topic 0--3, scenario 0--3, evidence 0--2) and keeps a document only if at least
one of its segments reaches the acceptance threshold on that rubric's $[0,1]$
projection: 0.5 in general, 0.3 for enumeration questions, corresponding to the
rubric's direct-answer band (Section~\ref{sec:gate}). It processes at most six
references per call, takes the
matched text from the server side rather than from the caller, and fails closed.
No surviving evidence means the caller must answer that the corpus does not
contain the answer.

\subsection{One query through the gate}
\label{sec:trace}

Q2 of Section~\ref{sec:eval-setup} is \emph{``EL inspection shows edge delamination:
root cause and corrective actions''}, asked through the public chat API with the
failure-analysis library pinned and the Hybrid lane selected.

\begin{enumerate}
\item \textbf{Policy.} The chat endpoint attaches the Hybrid lane's tool
whitelist server-side, so the agent cannot call a tool outside it.
\item \textbf{Candidate generation (one call).} The agent issues a single
retrieval call. In the shipped tool layer that call is a server-side
composition of \texttt{kb\_search\_vector}, the catalog tools \texttt{kb\_list}
and \texttt{kb\_get\_documents}, and the judging tool \texttt{kb\_laya\_judge};
\texttt{kb\_hybrid\_search} was an earlier-generation wrapper that is no longer
registered, so the configuration evaluated below was reproduced by composing
those tools. Two server-side channels then run: a dense vector channel
(the \texttt{kb\_search\_vector} call, top-30, cosine threshold 0.35) and a
catalog keyword channel over the pinned library's descriptions. Candidates
merge on (library, document), agreed documents rank first, and the pool is
capped at twelve.
\item \textbf{The gate.} The server fetches each candidate's text, segments it,
and scores \emph{does this text contain direct evidence that answers the query}
in batches of at most six. The agent supplies no text, so it cannot
argue for a document it has not read. On Q2 the gold report scores 0.846 while
five extruder-telemetry reports score 0.887--0.926.
\item \textbf{Cut and read budget.} A document-level relative cut keeps scores
$\geq \mathit{best}-0.15$ with a top-5 floor and a cross-domain anchor guard, so
an out-of-domain score cannot set the cut line; the top five survivors return in
full, at most 2\,000 characters each. That read window is the agent's only
full-content channel; the rest of the pool survives as scores in the audit
trail.
\item \textbf{The answer.} Citations name the document path, the lane, and the
engine score, graded P0 (direct evidence) through P2 (weak).
\item \textbf{The terminal paths.} If the cut leaves no survivor, Lane~A
escalates once to Lane~B; if Lane~B exhausts its budget, the run reports
not-found.
\end{enumerate}

Where the design fails, the same trace is explicit. Under the shipped re-ranking
role, Q2's gold (0.846) falls to rank eight behind the five higher-scoring
telemetry reports, outside the five-document read window, so the agent reads
confusers (Fig.~\ref{fig:anatomy}). Under the filter role the same pool, scores,
and verdicts leave the gold first: same gate, same evidence, only the role of the
induced ordering changes. This foreshadows the central result
(Section~\ref{sec:inflation}).

\subsection{The gate contract}
\label{sec:gate}

The gate is a locally hosted evidence-classification model, the Laya judging
service \cite{laya}, invoked through a \texttt{kb\_laya\_judge} tool. Given a
query and a batch of references, the server fetches and segments each document
and scores every segment under one of three criteria: \emph{evidence} (specific
substantiation; a topic match alone is insufficient), \emph{instance} (the text
contains a member of the requested class, routed automatically for enumeration
questions), or \emph{auto}. A document survives only if its best segment reaches
the threshold, so one at vector score 0.40 with a \emph{yes} verdict survives
while one at 0.95 that the judge rejects is dropped.

The score is documented as a 0--8 rubric (topic 0--3, scenario 0--3,
evidence 0--2; $\geq$6 answers directly, 5 escalates to the librarian lane,
$\leq$4 drops), with the $[0,1]$ score as its machine-readable projection. That
projection saturates (the logged maximum score is 1.0000), one reason thresholds
discriminate poorly between golds and same-domain confusers
(Section~\ref{sec:inflation}).

Every operating constant is a shipped production default: the thresholds, the
six-reference judge batch, the hybrid candidate cap of twelve, and the 0.15
relative-cut margin with its top-5 floor. They were fixed during earlier
deployments on a different corpus and were not tuned on the evaluation questions
of Section~\ref{sec:eval-setup}. The paper states their status rather than
selling the operating point: on the evaluated corpus none of the three
thresholds ever fired, because every candidate the retrievers returned already
scored far above them (Section~\ref{sec:inflation}). The gate's value there is
fail-closed provenance and the filter role, not discrimination; a deployment on
other content would have to re-tune these constants, and Section~\ref{sec:threats}
records that as a threat rather than a result. The classifier is a single, locally fine-tuned
evidence classifier whose checkpoint reference ships with the platform
artifacts; it is deterministic in our replays, and every verifier finding in this
paper inherits that choice, a scope limit noted in Section~\ref{sec:threats}.

Fail-closed semantics apply to every failure mode: fetch errors, engine errors,
and timeouts count as rejection. Survivors aggregate into an evidence pack with
per-passage provenance.

\subsection{Three lanes}

The three lanes are three points on a latency/coverage frontier
(Table~\ref{tab:lanes}; Fig.~\ref{fig:lanes}), not three implementations of one
idea.

\begin{table}[!htbp]
\centering
\scriptsize
\setlength{\tabcolsep}{3.5pt}
\caption{The three lanes as operating points. Measured figures are the
round-1 agentic runs of Section~\ref{sec:agentic}. All three share one
verification gate and one honesty contract; they differ only in retrieval
affordance and budget.}
\label{tab:lanes}
\begin{tabular}{p{1.85cm}>{\raggedright\arraybackslash}p{3.05cm}>{\raggedright\arraybackslash}p{2.15cm}>{\raggedright\arraybackslash}p{2.25cm}>{\raggedright\arraybackslash}p{3.15cm}}
\toprule
Lane & Optimises for & Tool budget & Measured (round 1) & Use when \\
\midrule
A. Search & latency on pinned libraries & 20-turn cap & 6--8 turns; 70--165\,s & interactive default; single-library questions \\
B. Librarian & exhaustive recall across all libraries & hard 12-call budget & 5--6 turns; 54--91\,s & must-not-miss; enumeration questions; no library pinned \\
C. Hybrid & single-call server-side orchestration & exactly 1 call & 3--5 turns; 36--70\,s & fastest grounded answer; latency-critical \\
\bottomrule
\end{tabular}
\end{table}

\begin{figure}[t]
\centering
\includegraphics[width=\textwidth]{fig2_lanes}
\caption{The three retrieval lanes. Every lane ends at the same
content-verification gate and the same honesty contract: zero verified
evidence escalates (Lane A$\to$B) or terminates in an explicit not-found
report.}
\label{fig:lanes}
\end{figure}

\textbf{Lane A (Search).} Phase A0 rewrites the query into (subject, attribute,
constraints) form, splits comparison questions into parallel sub-queries, and
routes incident-style queries to the experience store
(Section~\ref{sec:org-experience}) first. Phase A1 issues one wide dense query (top-30,
threshold 0.35, cross-library balancing), deduplicates by document, and keeps
candidate documents rather than truncating to top chunks. Phase A2 submits
citation-only references (library id plus document path) to the gate, and Phase
A3 returns the judge's evidence pack, this lane's only full-content channel, from
which the agent composes an answer whose citations carry confidence levels
(P0 direct evidence, P1 partial, P2 weak).

\textbf{Lane B (Librarian).} A fixed spine under a hard budget of twelve tool
calls. L0 reads the library catalog (ids, names, descriptions, counts) and prunes
shelves by the rewritten query; L1 lists each shelf's documents with
descriptions; L2 routes by query shape, sending literal and enumeration questions
(``which documents mention X'') to a deterministic scanner
(\texttt{scripts/enumerate\_scan.py}, which pulls the shelf's documents through
the tool layer and matches the requested terms literally and by regular
expression: $\approx$0.5\,s per document, zero LLM turns) and semantic questions
to judging; L3 judges in batches of at most six,
because larger batches historically triggered transport timeouts
and fail-closed batch loss; L4 reads only survivors, bounded to 2\,000 characters
and starting at the judge-reported evidence line; L5 emits a five-part report
(search paths, answer, sources, confidence, blind spots). The query-shape
choice is measured, not stylistic: on a probe with nine documented mentions, the
semantic judge recalled three of nine under the instance criterion at both
thresholds, while the scanner recalled nine of nine in 11.3\,s. Two rules are absolute:
catalog descriptions never prune documents by themselves, and an exhausted budget
terminates in blind-spot disclosure, not in a guessed answer.

\textbf{Lane C (Hybrid).} Agentic latency is dominated by LLM turns, not tools
(Section~\ref{sec:composition}), so Lane C spends exactly one call by moving the
pipeline into the server. A single call executes P1 (a
parallel block of a vector channel, top-30 at threshold 0.35, and a catalog
keyword channel); in the shipped tool layer that call is the composition of
\texttt{kb\_search\_vector}, \texttt{kb\_list}, \texttt{kb\_get\_documents} and
\texttt{kb\_laya\_judge} described above, not a single wrapper tool. P3 merges
candidates by (library, document), agreed-by-both
documents first, capped at twelve; P4 judges every candidate with the gate
(criterion \emph{evidence}, threshold 0.5); P4b applies the document-level
relative cut; and P5 reads the top five documents in full ($\leq$2\,000
characters each). The complete audit trail ships in the response, so a one-call
lane is still an inspectable lane.

\subsection{Selecting a lane is a budget decision}

The three lanes trade latency for recall under an enforced call and turn budget,
so the choice is how much budget to spend before the gate runs, not which
retriever is better.

What makes the lanes interchangeable is the gate. All three terminate in the same
contract: one gate, one fail-closed semantics, one evidence pack. Replacing a
lane changes the retrieval affordance and the budget, not the meaning of a
citation. Escalation follows a trigger, not a hunch: zero verified evidence
escalates Lane~A to Lane~B exactly once. Query-shape auto-selection is future
work (Section~\ref{sec:threats}).

One caveat is honest to record. On the evaluated corpus Lane~A occupies no
\emph{distinct} operating point, because Hybrid is faster and Librarian more
exhaustive. A two-lane deployment (Hybrid plus Librarian) is therefore defensible
today, with Search retained as the pinned-library interactive default; because
the lanes share one gate and one abstention semantics, that reduction costs
coverage, not trust.

\subsection{The two-stage entry point}
\label{sec:twostage}

One configuration scored in Section~\ref{sec:e4} is not one of the three lanes:
\path{kb_search_two_stage(query, kb_id, stage1_top_k=20, stage2_top_k=5, enable_graph_expansion=True, score_threshold=0.0, balance_kbs=False, slim=False)}.
It serves callers that want recall without paying for a verification lane, and
the platform's own retrieval skill names it the preferred entry point for agent
searches. Stage~1 builds a candidate document pool from the lexical index and,
when graph expansion is on, from the graph neighbours of the matched nodes
($\leq$20 documents); stage~2 runs the dense retriever inside that pool and
returns up to five chunks per document. A cross-library search that returns
documents from fewer than two libraries triggers an automatic supplementary
vector pass---the same lexical blind spot the librarian lane corrects by
escalating. It is fast because it never calls a judge: Hit@5 is 1.00 at 77\,ms
mean (Table~\ref{tab:retrieval}). That is also its limit, and the reason it is
not a lane: stage~2 returns chunks, so a caller that stops here holds candidate
text with no verdict attached, and any claim built on it is unverified by
construction. Only the gate turns a candidate into evidence.

\subsection{The gate as a platform primitive}
\label{sec:gate-primitive}

The gate is a server-side tool rather than a prompt:
\path{kb_laya_judge(query, documents, criterion=auto, threshold=0.5, max_evidence_chars=40000)}. The platform, not the agent, chooses the criterion
and fixes the threshold. Calls are serialized through a single judge worker with
a 900\,s call timeout, because an early concurrent workload crashed the worker
with a torch paging-file error; every verdict is logged with its query,
references, scores, and decision.

That placement makes verification a layer property rather than an agent
aspiration: any lane, harness, or external client inherits the gate, and none can
skip or re-threshold it per query. Enforcement has a second half in the lane
policy, a server-side per-lane tool whitelist: a librarian-lane agent is denied
the vector and two-stage tools and cannot reach around the catalog-first path it
was given. Denial is a property of that lane, not of the platform: the same
two-stage tool is scored directly as its own configuration in
Section~\ref{sec:e4}, and the vector tool backs Lane~A.

Externalized verification also distinguishes this design from verifiers inside
the generator. Self-RAG and CRAG verify during generation with the generator
itself \cite{selfrag,crag}; the gate sees only (query, document text) pairs and
is consumed as a tool. It is therefore reusable and inspectable, and equally
fallible: the LLM-judge critique applies here as much as anywhere
\cite{zheng-llmjudge}, which is why verdicts fail closed, are logged, and are
reported as evidence rather than as truth (Section~\ref{sec:inflation}). Deployment
cost is paid once: a twelve-candidate verification adds $\approx$0.5\,s to a call
(Table~\ref{tab:retrieval}), and a cold judge load costs 10--12\,s
(Section~\ref{sec:impl}).

\subsection{Honest abstention}

A threshold that never rejects a candidate cannot refuse a query. On four
out-of-corpus probes the judged hybrid still keeps 9--11 of its 12 candidates
(Section~\ref{sec:abstention}, Fig.~\ref{fig:ood}), and the graded probes behave
no differently: prompts with partial coverage and prompts with none both keep
10--12 of 12 candidates, with top scores up to 0.983
(Table~\ref{tab:robustness} and Section~\ref{sec:e7}). Because the $[0,1]$
projection saturates, no cut point separates a supported answer from an
unsupported one.

Abstention therefore lives in the reading-and-answering layer, enforced by the
lane contract at answer time rather than advised in a prompt. A caller with no
surviving evidence must report a not-found with its search scope: which libraries
and how many documents were searched, how the query was routed, and which topics
sit outside the corpus. The escalation rule makes that state terminal, so an
exhausted budget produces a report and not a guess.

Abstention is thus a terminal state evaluated like any other outcome: all six
out-of-corpus probe runs returned explicit not-found reports with zero fabricated
citations, even though the tool layer had handed them candidates
(Sections~\ref{sec:agentic} and \ref{sec:abstention}). The division of labour is
the point: the gate decides what may be cited, the reading layer decides whether
to answer at all, and the measurements put the reliability in the second and not
the first.
